\chapter{Conclusion and Future Outlook}
\label{ch:conclusion_and_future_outlook}

% Core Research Question:
% What general principles have been established, what are the fundamental limitations
% of this study, and what is required to translate data-driven closed-loop neurostimulation
% toward clinical reality?

% Contents:
% - Summary of Key Contributions and Findings:
%   * Feasibility of non-invasive, data-driven closed-loop MPC for seizure suppression without internal state access.
%   * Superiority of spectral Observables over raw time-domain waveforms for neural mass dynamics.
%   * Effectiveness of constrained receding-horizon current optimization under strict Kirchhoff safety budgets.
% - Critical Discussion of Limitations:
%   * Model fidelity: Lumped three-population Jansen-Rit dynamics vs. microscopic ion channels.
%   * Electrode geometry: Fixed 3-electrode montage
%   * Biological adaptation: Absence of long-term synaptic plasticity, kindling, or tissue habituation.
% - Future Directions and Translational Roadmap:
%   * Online adaptive model updates (closing the loop on non-stationary brain dynamics).
%   * Robust and stochastic MPC formulations to handle tissue conductivity and leadfield uncertainties.
%   * Experimental roadmap from in-silico whole-brain models to animal preparations and clinical translation.
